% Part: turing-machines
% Chapter: machines-computations
% Section: introduction

\documentclass[../../../include/open-logic-section]{subfiles}

\begin{document}

\olfileid{tur}{mac}{int}
\olsection{ভূমিকা}

ধরা যাক, $\Nat$ থেকে $\Nat$-এ একটি অপেক্ষক; তার \emph{গণনসাধ্য}
হওয়ার অর্থ কী? প্রথম দিকের উত্তরগুলির মধ্যে সবচেয়ে পরিচিত একটি হলো:
কোনো টুরিং যন্ত্র দিয়ে অপেক্ষকটিকে গণনা করা গেলে সেটি গণনসাধ্য। অ্যালান
টুরিং ১৯৩৬ সালে এই ধারণা উপস্থাপন করেন। টুরিং যন্ত্র একটি
\emph{গণনামূলক মডেল}—“গণন-পদ্ধতি” ধারণাটিকে গাণিতিকভাবে নির্ভুল করে
তোলার একটি উপায়। কথাটির সঠিক অর্থ কী, তা নিয়ে বিতর্ক আছে; তবে টুরিং
যন্ত্র যে গণন-পদ্ধতি নির্দিষ্ট করার একটি উপায়, সে বিষয়ে ব্যাপক ঐকমত্য
আছে। “টুরিং যন্ত্র” কথাটি চলমান যন্ত্রাংশ-সহ কোনো ভৌত যন্ত্রের ছবি মনে
আনলেও, কঠোর অর্থে এটি নিছক একটি গাণিতিক নির্মাণ; সেই হিসেবে এটি
গণন-পদ্ধতির ধারণাকে আদর্শায়িত করে। যেমন, টুরিং যন্ত্রের সময় বা স্মৃতির
চাহিদার উপর আমরা কোনো সীমা দিই না: কোনো গণনায় মহাবিশ্বে যত পরমাণু আছে
তার চেয়েও বেশি সঞ্চয়স্থান বা বেশি ধাপ লাগলেও টুরিং যন্ত্র সেটি গণনা
করতে পারে।

\begin{explain}
টুরিং যন্ত্রকে বিশেষ ধরনের একটি কাল্পনিক ব্যবস্থার প্রোগ্রাম হিসেবে ভাবাই
সম্ভবত সবচেয়ে ভালো। এই ব্যবস্থায় একটি \emph{ফিতা} এবং একটি
\emph{পাঠ-লেখ হেড} থাকে। টুরিং যন্ত্রের যে সংস্করণ আমরা ব্যবহার করব,
তাতে ফিতাটি এক দিকে (ডান দিকে) অসীম এবং কতগুলি \emph{ঘরে} বিভক্ত;
প্রতিটি ঘরে একটি সসীম \emph{বর্ণমালা} থেকে একটি প্রতীক থাকতে পারে। এমন
বর্ণমালায় যেকোনো সসীম সংখ্যক ভিন্ন প্রতীক থাকতে পারে, তবে আমাদের প্রধানত
তিনটিতেই কাজ চলবে: $\TMendtape$, $\TMblank$ এবং $\TMstroke$। ব্যবস্থাটি
চালু করার সময় ফিতাটি ফাঁকা থাকে (অর্থাৎ প্রতিটি ঘরে $\TMblank$ প্রতীক
থাকে), শুধু বাঁদিকের প্রথম ঘরে থাকে $\TMendtape$, আর সসীমসংখ্যক ঘরে থাকে
\emph{নিবেশ}। যেকোনো সময় ব্যবস্থা সসীমসংখ্যক \emph{দশা}-র একটিতে থাকে।
শুরুতে হেডটি নিবেশের জন্য নির্ধারিত বাঁদিকের প্রথম ঘরটি পড়ে এবং একটি
নির্দিষ্ট \emph{আরম্ভিক দশা}-য় থাকে। ব্যবস্থাটির চালনার প্রতিটি ধাপে
হেডটি তখন যে ঘর পড়ছে তার বিষয়বস্তু এবং ব্যবস্থার বর্তমান দশা, টুরিং
যন্ত্রের প্রোগ্রামের সঙ্গে মিলে, পরের ঘটনাটি নির্ধারণ করে। টুরিং-যন্ত্রের
প্রোগ্রামটি একটি আংশিক অপেক্ষক: সেটি নিবেশ হিসেবে একটি দশা~$q$ ও একটি
প্রতীক~$\sigma$ নেয় এবং একটি ত্রয়ী~$\tuple{q', \sigma', D}$ নির্গত
করে। ব্যবস্থা যখন দশা~$q$-তে থেকে প্রতীক~$\sigma$ পড়ে, তখন বর্তমান
ঘরের প্রতীকটির জায়গায় $\sigma'$ লেখে; $D$ যথাক্রমে $\TMleft$,
$\TMright$ বা $\TMstay$ হলে হেডটি বাঁ দিকে যায়, ডান দিকে যায়, অথবা
একই জায়গায় থাকে; এবং ব্যবস্থা দশা~$q'$-তে যায়।
% BN-SRC-162 TARGET-CORRECTION-BEGIN
\par\smallskip\noindent\textnormal{\small\textbf{উৎস-সংশোধন (BN-SRC-162):} হিমায়িত ইংরেজি ভূমিকায় শুরুতে হেডটি ফিতার একেবারে বাঁদিকের ঘর, অর্থাৎ $\TMendtape$-চিহ্নিত ঘরটি পড়ে বলা হয়েছে। এই অধ্যায়ের আনুষ্ঠানিক আরম্ভিক কনফিগারেশন $\tuple{\TMendtape \frown I, 1, q_0}$ অনুযায়ী হেডটি তার ডান পাশের, নিবেশের জন্য নির্ধারিত প্রথম ঘরটি পড়ে। বাংলা পাঠে সেই অবস্থান রাখা হয়েছে।} % BN-SRC-162 TARGET-CORRECTION-END

যেমন, \olref{fig:tm}-এর পরিস্থিতিটি বিবেচনা করো।
\begin{figure}
  \olasset{\olpath/assets/diagrams/turing-machine.tikz}
  \caption{একটি টুরিং যন্ত্র তার প্রোগ্রাম চালাচ্ছে।}
  \ollabel{fig:tm}
\end{figure}
টুরিং যন্ত্রটির ফিতার দৃশ্যমান অংশে বাঁদিকের প্রথম ঘরে ফিতার শেষ-প্রতীক
$\TMendtape$, তার পরে তিনটি $1$, একটি $0$, এবং আরও চারটি $1$ আছে। হেডটি
বাঁ দিক থেকে তৃতীয় ঘরটি পড়ছে, সেখানে একটি~$1$ আছে, আর যন্ত্রটি
দশা~$q_1$-এ রয়েছে—আমরা বলি, “যন্ত্রটি দশা~$q_1$-এ একটি $1$ পড়ছে।”
টুরিং-যন্ত্রের প্রোগ্রামটি নিবেশ $\tuple{q_1, 1}$-এর জন্য ত্রয়ী
$\tuple{q_2, 0, \TMstay}$ ফেরত দিলে যন্ত্রটি তৃতীয় ঘরের~$1$-এর জায়গায়
একটি~$0$ লিখত, পাঠ-লেখ হেডটিকে সেখানেই রাখত এবং দশা~$q_2$-এ চলে যেত।
তারপর প্রোগ্রামটি নিবেশ $\tuple{q_2, 0}$-এর জন্য
$\tuple{q_3, 0, \TMright}$ ফেরত দিলে যন্ত্রটি~$0$-টির উপর আরেকটি~$0$
লিখত (ফলে কার্যত পাঠ-লেখ হেডের নিচে ফিতার বিষয়বস্তু অপরিবর্তিত থাকত),
এক ঘর ডান দিকে যেত এবং দশা~$q_3$-এ প্রবেশ করত। এভাবেই চলতে থাকত।

আমরা বলি, যন্ত্রটি \emph{থামে}, যখন সেটি এমন কোনো দশা~$q_n$ ও
প্রতীক~$\sigma$-র সম্মুখীন হয়, যার জন্য $\tuple{q_n, \sigma}$-এর কোনো
নির্দেশ নেই; অর্থাৎ নিবেশ $\tuple{q_n,\sigma}$-এর জন্য অবস্থান্তর অপেক্ষক
অসংজ্ঞায়িত। অন্যভাবে বললে, যন্ত্রটির পালন করার মতো কোনো নির্দেশ থাকে না
এবং তখন সেটি কাজ বন্ধ করে। কখনও একটি বিশেষ থামা-দশা~$h$ দিয়ে থামা
প্রকাশ করা হয়। পরে এটি আরও বিশদে দেখানো হবে।
\end{explain}

\begin{digress}
টুরিংয়ের “On computable numbers” গবেষণাপত্রটির সৌন্দর্য হলো, সেখানে তিনি
শুধু একটি আনুষ্ঠানিক সংজ্ঞাই দেননি; সংজ্ঞাটি গণনসাধ্যতার স্বজ্ঞাত ধারণাকে
ধারণ করে—এই যুক্তিও দিয়েছেন। সংজ্ঞা থেকেই পরিষ্কার হওয়া উচিত যে টুরিং
যন্ত্রে গণনসাধ্য যেকোনো অপেক্ষক স্বজ্ঞাত অর্থে গণনসাধ্য। বিপরীত কথাটির
পক্ষে—অর্থাৎ স্বাভাবিকভাবে গণনসাধ্য বলে আমরা যে অপেক্ষককে গণ্য করব, সেটি
এমন কোনো যন্ত্রে গণনসাধ্য—টুরিং তিন ধরনের যুক্তি দেন। তাঁর ভাষায় সেগুলি:
\begin{enumerate}
\item স্বজ্ঞার কাছে সরাসরি আবেদন।
\item দুটি সংজ্ঞার সমতুল্যতার প্রমাণ (নতুন সংজ্ঞাটি যদি স্বজ্ঞার কাছে
  বেশি গ্রহণযোগ্য হয়)।
\item গণনসাধ্য সংখ্যার বড় বড় শ্রেণির উদাহরণ দেওয়া।
\end{enumerate}
আমাদের লক্ষ্য “নীতিগতভাবে” গণনসাধ্যতার ধারণাকে সংজ্ঞায়িত করা; অর্থাৎ
সময় ও স্থানের ব্যবহারিক সীমাবদ্ধতা হিসাব না করা। অবশ্য গণনসাধ্যতার
সবচেয়ে ব্যাপক সংজ্ঞাটি স্থির হয়ে গেলে সীমিত সম্পদে গণনা নিয়ে ভাবা যায়;
সেই আলোচনাই “গণনামূলক জটিলতা” নামে পরিচিত বিষয়টির কেন্দ্র।
\end{digress}

\begin{history}
অ্যালান টুরিং ১৯৩৬ সালে টুরিং যন্ত্র উদ্ভাবন করেন। তখন তাঁর আগ্রহ ছিল
প্রথম-ক্রমের যুক্তিবিদ্যার নির্ণেয়তা; তবু গবেষণাপত্রটিকে কম্পিউটার নকশার
ভিত্তি বিষয়ে একটি নির্ধারক রচনা বলা হয়েছে। সেখানে টুরিং গণনসাধ্য বাস্তব
সংখ্যার উপর মন দেন, অর্থাৎ যেসব বাস্তব সংখ্যার দশমিক বিস্তার গণনা করা
যায়; তবে তিনি লক্ষ করেন, স্বাভাবিক সংখ্যার উপর গণনসাধ্য অপেক্ষক ইত্যাদির
ক্ষেত্রে তাঁর ধারণাগুলি খাপ খাওয়ানো কঠিন নয়। খেয়াল করো, প্রথম সচল
সাধারণ-উদ্দেশ্য কম্পিউটারটি ১৯৪১ সালে নির্মিত হওয়ার (জার্মান কনরাড সুসে
তাঁর মা-বাবার বসার ঘরে সেটি বানিয়েছিলেন) পুরো পাঁচ বছর আগে এই কাজ;
ব্লেচলি পার্কে টুরিং ও তাঁর সহকর্মীরা সংকেতভেদী কলসাস বানানোর (১৯৪৩)
সাত বছর আগে; মার্কিন এনিয়াকের (১৯৪৫) নয় বছর আগে; ম্যানচেস্টারে প্রথম
ব্রিটিশ সাধারণ-উদ্দেশ্য কম্পিউটার—ম্যানচেস্টার স্মল-স্কেল এক্সপেরিমেন্টাল
মেশিন—নির্মিত হওয়ার (১৯৪৮) বারো বছর আগে; এবং মার্কিনদের বিন্যাক প্রথম
পরীক্ষার (১৯৪৯) তেরো বছর আগে। ম্যানচেস্টার SSEM-ই প্রথম সঞ্চিত-প্রোগ্রাম
কম্পিউটার হওয়ার স্বাতন্ত্র্য বহন করে—এর আগের যন্ত্রগুলিকে প্রতিটি নতুন
কাজের জন্য হাতে তার জুড়ে নতুনভাবে সাজাতে হতো।
\end{history}

\end{document}
